% R585: EGA I, 0, 3.1.1--3.1.3, con ambas notas.
% Traducción: OpenAI Codex, GPT-6.1 Sol, esfuerzo Ultra.
\section{Complementos sobre los haces}
\label{section:0.3-fr}

\subsection{Haces con valores en una categoría}
\label{subsection:0.3.1-fr}

\begin{env}[3.1.1]
\label{0.3.1.1-fr}
Sean $\boldsymbol{K}$ una categoría, $(A_\alpha)_{\alpha\in I}$,
$(A_{\alpha\beta})_{(\alpha,\beta)\in I\times I}$ dos familias de objetos de
$\boldsymbol{K}$ tales que $A_{\beta\alpha}=A_{\alpha\beta}$, y
$(\rho_{\alpha\beta})_{(\alpha,\beta)\in I\times I}$ una familia de morfismos
$\rho_{\alpha\beta}:A_\alpha\to A_{\alpha\beta}$. Diremos que un par
formado por un objeto $A$ de $\boldsymbol{K}$ y una familia de morfismos
$\rho_\alpha:A\to A_\alpha$ es \emph{solución del problema universal} definido
por los datos de las familias $(A_\alpha)$, $(A_{\alpha\beta})$ y
$(\rho_{\alpha\beta})$ si, para todo objeto $B$ de $\boldsymbol{K}$, la aplicación
que a todo $f\in\operatorname{Hom}(B,A)$ hace corresponder la familia
$(\rho_\alpha\circ f)\in\prod_\alpha\operatorname{Hom}(B,A_\alpha)$ es una
\emph{biyección} de $\operatorname{Hom}(B,A)$ sobre el conjunto de las $(f_\alpha)$
tales que
$\rho_{\alpha\beta}\circ f_\alpha=\rho_{\beta\alpha}\circ f_\beta$ para todo
par de índices $(\alpha,\beta)$. Se ve inmediatamente que, si existe tal
solución, es única salvo isomorfismo.
\end{env}

\begin{env}[3.1.2]
\label{0.3.1.2-fr}
No recordaremos la definición de un \emph{prehaz}
$U\to\mathcal{F}(U)$ sobre un espacio topológico $X$, con valores en una
categoría $\boldsymbol{K}$ (G, I, 1.9); diremos que tal prehaz es un
\emph{haz con valores en $\boldsymbol{K}$} si satisface el axioma
siguiente:

\smallskip
\noindent
\textup{(F)} \emph{Para todo recubrimiento $(U_\alpha)$ de un abierto $U$ de $X$
por abiertos $U_\alpha$ contenidos en $U$, si se denota por $\rho_\alpha$
\textup{(resp. $\rho_{\alpha\beta}$)} el morfismo de restricción}
\[
  \mathcal{F}(U)\to\mathcal{F}(U_\alpha)
  \qquad
  \bigl(\text{resp. }\mathcal{F}(U_\alpha)
  \to\mathcal{F}(U_\alpha\cap U_\beta)\bigr),
\]

\oldpage[0\textsubscript{I}]{24}

\emph{el par formado por $\mathcal{F}(U)$ y la familia $(\rho_\alpha)$ es
solución del problema universal para $(\mathcal{F}(U_\alpha))$,
$(\mathcal{F}(U_\alpha\cap U_\beta))$ y $(\rho_{\alpha\beta})$ \textup{(3.1.1)}}%
\footnote{Es un caso particular de la noción general de \emph{límite
proyectivo} (no filtrante) (\emph{véase} (T, I, 1.8) y el libro en preparación
anunciado en la Introducción).}.

Equivale a decir que, para todo objeto $T$ de $\boldsymbol{K}$, la
familia $U\to\operatorname{Hom}(T,\mathcal{F}(U))$ es un \emph{haz
de conjuntos}.
\end{env}

\begin{env}[3.1.3]
\label{0.3.1.3-fr}
Supongamos que $\boldsymbol{K}$ es la categoría definida por una «especie de
estructura con morfismos» $\Sigma$, siendo los objetos de $\boldsymbol{K}$, por tanto,
los conjuntos dotados de estructuras de especie $\Sigma$ y los morfismos los de
$\Sigma$. Supongamos que la categoría $\boldsymbol{K}$ satisface además la
condición siguiente:

\smallskip
\noindent
\textup{(E)} Si $(A,(\rho_\alpha))$ es solución de un problema de aplicación
universal \emph{en la categoría $\boldsymbol{K}$} para unas familias
$(A_\alpha)$, $(A_{\alpha\beta})$, $(\rho_{\alpha\beta})$, entonces es también
solución del problema de aplicación universal para las mismas familias
\emph{en la categoría de conjuntos} (es decir, cuando se consideran $A$,
los $A_\alpha$ y $A_{\alpha\beta}$ como conjuntos, los $\rho_\alpha$ y
$\rho_{\alpha\beta}$ como aplicaciones)%
\footnote{Se puede demostrar que esto significa también que el funtor canónico
$\boldsymbol{K}\to(\mathrm{Ens})$ \emph{conmuta con los límites proyectivos} (no
necesariamente filtrantes).}.

En estas condiciones, la condición (F) implica que, considerado como
prehaz de \emph{conjuntos}, $U\to\mathcal{F}(U)$ es un \emph{haz}.
Además, para que una aplicación $u:T\to\mathcal{F}(U)$ sea un morfismo de
$\boldsymbol{K}$, es necesario y suficiente, en virtud de (F), que cada aplicación
$\rho_\alpha\circ u$ sea un morfismo $T\to\mathcal{F}(U_\alpha)$, lo que
significa que la estructura de especie $\Sigma$ sobre $\mathcal{F}(U)$ es
\emph{estructura inicial} para los morfismos $\rho_\alpha$. Recíprocamente,
supongamos que un prehaz $U\to\mathcal{F}(U)$ sobre $X$, con valores en
$\boldsymbol{K}$, es un \emph{haz de conjuntos} y satisface la condición
anterior; es claro entonces que satisface (F), por lo que es un
\emph{haz con valores en $\boldsymbol{K}$}.
\end{env}
